\documentclass[two column]{article}
\usepackage{graphicx} 
\usepackage[letterpaper,top=2cm,bottom=2cm,left=3cm,right=3cm,marginparwidth=1cm]{geometry}
\usepackage[colorlinks=true, allcolors=blue]{hyperref}

\title{Smart Wearable System for Pediatric Asthma Monitoring}
\author{Bozhena Bursalı, Selin Atakan}
\date{April 2026}

\begin{document}
\maketitle
\begin{abstract}
Managing pediatric asthma is difficult because most current tools are reactive, only acting after symptoms appear. This paper introduces an IoT-based wearable system that monitors both a child's health (pulse and $SpO_2$) and their environment (air quality and humidity) in real-time. Using an ESP32-based prototype and data fusion, the system aims to predict asthma attacks before they happen. Our goal is to provide a low-cost, proactive solution that helps families manage asthma more effectively and reduces emergency room visits.
\end{abstract}

\section{Introduction}
Asthma is a major health issue for children, often triggered by environmental factors like air pollution and humidity[cite: 12]. Currently, most families manage asthma "reactively," meaning they only use medicine after a child starts struggling to breathe. This increases the risk of serious attacks and hospital visits[cite: 17].

While there are many asthma apps available, most lack real-time data. For example, apps like AsthmaMD require manual input, and systems like Propeller Health do not track heart rate or oxygen levels. There is a clear need for a tool that combines body sensors with environmental data.

In this project, I develope a smart wearable prototype using an ESP32 microcontroller[cite: 7.2]. Our system collects physiological data (pulse, $SpO_2$) and environmental data (PM2.5, temperature) to provide an early warning for potential attacks[cite: 12]. By using "data fusion," we can offer a more accurate and proactive way to keep children safe.\href{https://www.who.int/news-room/fact-sheets/detail/asthma}{World Health Organization}

\subsection{}



\end{document}